\section{Dataset Preparation}
\label{sec:datasets}
\suppressfloats[t]
This section documents the completed model fixture and graph cohorts, defines identity labels, and specifies the local verification inputs for the seven planned experiments in RQ1--RQ3. The primary recovery evaluation takes supplied JavaScript and fixed, versioned references; automatic library identification and reference retrieval are input preparation, not evaluated contributions. Source maps are withheld from model-driven recovery. The separate graph simulations use an explicitly declared oracle. The practical-use dataset for RQ4 remains to be defined.

\subsection{Completed model fixture}
\label{sec:model-fixture}
The model pilot uses a crafted task-board application and the published jQuery~4.0.0 ESM distribution. Three evaluated browser builds use esbuild~0.25.9, Rollup~4.50.0, and webpack~5.97.1 with minification enabled and tree shaking disabled. The distribution contains 558 reference definitions, and the three builds provide 1629 supported jQuery occurrences. These inputs were used during method development and are not an unseen-repository evaluation. Appendix~\ref{sec:pilot-setup} preserves build settings, model deployments, and the original experimental procedure.

\subsection{Completed graph cohorts and bundler coverage}
\label{sec:bundler-selection}
The completed oracle-assisted graph experiment extends beyond the crafted jQuery fixture. The selection is motivated by the State of JavaScript 2025 build-tool usage ranking~\cite{stateofjs2025}. We prioritize the dedicated bundler entries: webpack, esbuild, Rollup, Turbopack, Parcel, Rolldown, and Rspack (Table~\ref{tab:bundler-usage}). The survey also lists compilers, wrappers, and runtimes; we do not count them as additional dedicated bundler families. Vite's 84\% published usage motivates coverage of its backends: earlier Vite versions combine esbuild and Rollup, whereas Vite~8 adopts Rolldown~\cite{vite8}. This is a rationale for protocol coverage, not an estimate of production deployment share or a probability sample of applications.

\begin{table}[t]
\caption{Dedicated bundler entries ordered by published State of JavaScript 2025 usage ratios~\cite{stateofjs2025}. Percentages retain the survey API\textquotesingle s rounded ratios. Usage means having used a tool, not current production market share.}
\label{tab:bundler-usage}
\centering\small
\begin{tabular}{lr@{\qquad}lr}
\toprule
Bundler & Usage (\%) & Bundler & Usage (\%)\\
\midrule
\nolinkurl{webpack} & 87 & \nolinkurl{parcel} & 26\\
\nolinkurl{esbuild} & 54 & \nolinkurl{rolldown} & 10\\
\nolinkurl{rollup} & 47 & \nolinkurl{rspack} & 7\\
\nolinkurl{turbopack} & 28 &  & \\
\bottomrule
\end{tabular}
\end{table}


We retain three distinct cohorts. An earlier esbuild graph study evaluates three pinned Angular applications and 25 repository/library cases (7497 supported pairs). The current analyzer's development cohort contains eight Rspack, Rolldown, Turbopack, and Parcel applications, with 69 cases (62 having supported targets) and 161,329 supported pairs. These repositories were used while developing the adapters. The new-repository cohort freezes the analyzer before evaluation and selects five previously untested repositories: \nolinkurl{ionic-team/ionic-conference-app}, \nolinkurl{geeker-admin/Geeker-Admin}, \nolinkurl{robustmq/robustmq-copilot}, \nolinkurl{decolua/9router}, and \nolinkurl{cobiwave/simplefolio}. It covers one repository per non-webpack family, with 41 cases (40 having supported targets) and 28,780 supported pairs. The cohorts are reported separately, and the earlier esbuild implementation is not pooled with the current revision. Rollup and webpack additionally occur in the controlled jQuery builds; they are not new-repository subjects in this experiment.

We select at most ten loaded packages per application by deduplicated, source-map-embedded JavaScript-family source bytes, before inspecting recovery outcomes. The established exclusions remove the selected Angular, React, and Vue framework runtimes and build loaders. Component libraries such as Ionic and Element Plus remain in scope; this operational boundary is not a universal framework/library taxonomy. Only emitted packages count: a dependency declared but absent from bundled JavaScript does not become a test library. Simplefolio consequently contributes only vanilla-tilt. Artifacts include supplied lazy, shared, and runtime chunks within the selected browser or renderer build, rather than only the files reached by a navigation trace. Separate service-worker scripts are excluded from the new-repository browser scope.

Every subject retains a pinned commit, exact JS/map hashes, build configuration, selected packages, and input inventory. Existing corpus artifacts are reused where available. Geeker Admin is built with native Vite~8.0.16 and source maps; RobustMQ uses native Rsbuild~1.4.1 with a configuration wrapper enabling JavaScript source maps; Simplefolio uses native Parcel~2.8.2. For 9router, the pinned application script uses webpack, so we explicitly rebuild it with Next~16.3.5's Turbopack backend. Its existing build configuration emits browser maps. Google Fonts fetching failed; the successful isolated offline build uses Next's mock-response hook to substitute an empty CSS response for the Inter font URL. This omits font assets without editing application JavaScript. Failed attempts, the fixture, and the effective configuration are preserved. We therefore describe this subject as a build-adjusted Turbopack experiment, rather than an untouched default release.

These completed graph cohorts use source-map-selected references and supplied emitted artifacts. Their oracle-assisted design is specified in Section~\ref{sec:connectivity-design}. The development cohort cannot be relabeled as unseen; only the five-repository graph check freezes the analyzer before subject evaluation. None of these graph cohorts constitutes a completed model-driven recovery evaluation on unseen repositories.

\subsection{Identity labels and measurement units}
\label{sec:identity-measurements}
The following source-map policy defines assessable identities in the completed experiments and will be retained for the planned verification evaluation. In model-driven runs, the oracle labels accepted matches for scoring; it does not supply candidates or correct model decisions.

For each mapping segment, we locate the innermost bundled function containing its generated position. We retain only positions inside that function's own body, excluding the function header and the opening and closing brace positions. Expression-bodied arrows retain the full expression interval. Nested functions receive their own mapping points rather than contributing them to an enclosing function. Original positions then vote for the innermost reference function containing them.

A bundled function receives a supported reference identity when at least two retained points vote for the same reference function, its votes account for at least 80\% of all retained mapped points, and it is the unique highest-vote function. We also retain sufficiently supported application or bundler-runtime origins; predicting a library identity for one of these functions counts as wrong. A function without a supported identity or non-library origin remains unresolved. A predicted match on such a function is unassessed, rather than being silently counted as correct or wrong.

The oracle supports positional identity, not formal semantic equivalence. A different reference identity can still implement equivalent behavior. Conversely, a source-linked function can undergo a behavior-preserving specialization that changes its behavior outside the actual calling context. Our measured errors concern recovering reference identities under this oracle. Functions created by the build without a supported reference counterpart are outside the target population; their library authorship is not inferred from appearance.

Let $D$ denote supported occurrences of the selected library functions, and let $C$, $W$, and $U$ denote correct, wrong, and unassessed accepted matches. Recovery is $C/D$; assessed identity precision is $C/(C+W)$, with $U$ reported separately. Wrong library assignments to supported application or runtime functions contribute to $W$ even though those functions are outside $D$. We retain reference inventory size, parsing failures, and unsupported or ambiguous transformations so that exclusions remain visible. Distinct reference identities, bundled occurrences, library cases, repositories, and website findings have different denominators. Report per-library and per-bundler results, macro summaries, and pooled counts without treating repeated occurrences as independent trials.

\subsection{Planned verification dataset}
\label{sec:planned-verification}
The planned local dataset will cover different bundlers and target libraries/frameworks for experiments~1--5. Target packages, exact versions, framework-runtime inclusion, and supplied artifact scope will be selected before inspecting recovery outcomes. Existing exclusions in the completed graph cohorts remain unchanged and will not be treated as framework recovery results. Configuration selection for experiments~6--7 will use repository-disjoint development subjects, and the analyzer, models, prompts, and thresholds will be frozen before held-out testing. Dataset sizes, selected repositories, and preparation outcomes remain pending. Existing development subjects will retain their development status.

For each repository, we will preserve its commit, dependency lockfile, effective bundler configuration and version, browser or renderer build, and JS/map hashes. The declared local artifact set will identify entry, lazy, shared, and runtime chunks and their execution environments. All compared conditions within an experiment will use the same supplied files and local context. Locally emitted artifacts will not be described as browser-intercepted code or as complete runtime coverage. Source maps and manual review labels will be stored separately from recovery inputs.

Versioned references will be prepared from the repository's pinned dependencies and corresponding published distributions, with package variants, integrity metadata, and file hashes retained. The reference set will be fixed before recovery and shared across compared configurations. Missing or incompatible reference artifacts will be recorded as preparation exclusions, together with their affected packages; evaluation labels will not silently repair inputs. This design measures function recovery conditional on the declared references, not automatic package identification or retrieval.

For the planned boundary characterization in Section~\ref{sec:boundary-characterization}, we will retain contributor relations separately from the dominant-identity oracle above and include generated top-level regions when reviewing no-target candidates. Each source inventory will state whether it covers complete build inputs or only map-embedded source texts; unrepresented modules and unavailable source contents will limit the claimed population. No source-map absence alone will establish removal. The resulting manually validated counts will be linked to each library/build case's recovery measurements.

\subsection{Practical-use dataset: to be filled}
\label{sec:planned-websites}
The application task, dataset, and evaluation protocol for RQ4 remain to be filled. No website population, downstream comparison, or practical-use result is specified by the current roadmap.
